\documentclass[a4paper,twoside,twocolumn]{IEEEtran}
% Use one column for design and twocolumn for research reports


% --------------------------------------------------
% Packages
% --------------------------------------------------

\usepackage{graphicx}                       % Images
\usepackage[caption=false,font=footnotesize]{subfig} % Subfigures
\usepackage{amsmath,amsfonts,amssymb}       % Equations
\usepackage{algorithm}                      % Algorithms
\usepackage{algorithmic}
\usepackage{cite}                           % References
\usepackage{enumitem}                       % Lists
\usepackage{listings}                       % Code blocks
\usepackage{xcolor}                         % Colours for code highlighting
\usepackage{multirow}                       % Cells spanning multiple rows
\usepackage{array}                          % Extra column types for tables
\usepackage{booktabs}                       % Professional table rules
\usepackage{url}                            % Typesetting web addresses


% You can define unsual hypenation here.
\hyphenation{op-tical net-works semi-conduc-tor IEEE-Xplore}


% --------------------------------------------------
% Code listing style (used by all lstlisting blocks)
% Change `language=` to Java, Python, etc. as needed.
% --------------------------------------------------

\definecolor{codegreen}{rgb}{0.0,0.5,0.0}
\definecolor{codeblue}{rgb}{0.0,0.0,0.6}
\definecolor{codegray}{rgb}{0.5,0.5,0.5}

\lstset{
    language=C++,
    linewidth=0.95\columnwidth,
    xleftmargin=0.05\columnwidth,
    basicstyle=\ttfamily\footnotesize,
    keywordstyle=\bfseries\color{codeblue},
    commentstyle=\itshape\color{codegreen},
    stringstyle=\color{codegreen},
    numbers=left,
    numberstyle=\tiny\color{codegray},
    numbersep=6pt,
    frame=single,
    breaklines=true,
    showstringspaces=false,
    captionpos=b,
    tabsize=2,
    columns=fixed
}


% --------------------------------------------------
% Document
% --------------------------------------------------

\begin{document}


% --------------------------------------------------
% Title
% --------------------------------------------------

\title{A Sample Project Template\\ for Spiltek Students}

\author{Elev En, Elev To, Elev Tre, Elev Fire, Elev Fem% <-this % stops a space
  % Here you can specify the roles and contributions of the team.
  \thanks{\textbf{Lead Designer:} Elev En \dotfill 20\%}%
  \thanks{\textbf{Systems Designer:} Elev To \dotfill 20\%}%
  \thanks{\textbf{Lead Programmer:} Elev Tre \dotfill 20\%}%
  \thanks{\textbf{QA Lead:} Elev Fire \dotfill 20\%}%
  \thanks{\textbf{Art Director:} Elev Fem \dotfill 20\%}
}

% The paper headers
\markboth{Syddansk Universitet (SDU) - Spiludvikling og Læringsteknologi}%
% Here mark the surname of the authors and the tile again.
{En \MakeLowercase{\textit{et al.}}: A Sample Project Template for Spiltek Students}

\maketitle


% --------------------------------------------------
% Abstract
% --------------------------------------------------

\begin{abstract}
This document is a template for Spiltek project reports. Your abstract should
briefly state the game or system you built, the approach you took, and your main
result --- for example, what a playtest revealed about your core mechanic.
\end{abstract}


\begin{IEEEkeywords}
game design, playtesting, prototyping, level design
\end{IEEEkeywords}


% --------------------------------------------------
% Introduction
% --------------------------------------------------

\section{Introduction}

LaTeX is a typesetting system in which you write plain text marked up with simple
commands, and it produces a polished, consistently formatted document. Its main
advantage is that you concentrate on \emph{what} you are writing while LaTeX
handles \emph{how} it looks: it manages spacing, fonts, section numbering, and
the placement of figures (see Fig.~\ref{fig:single}) and tables for you, 
so the layout takes care of itself and stays uniform throughout. 
The rest of this document introduces the pieces you will need for a Spiltek report.

This is a normal paragraph. In LaTeX you do not press Enter to start a new
paragraph. Instead, a paragraph continues until you leave a completely blank
line; a single line break in the source is treated as an ordinary space, so you
can wrap your source text however you like without affecting the output. One or
more blank lines end the current paragraph and begin the next one (extra blank
lines have no additional effect). Avoid using \verb|\\| to separate paragraphs:
that command forces a line break \emph{within} a paragraph and does not start a
new one, so it usually produces the wrong spacing.

This is the second paragraph. Notice the blank line above it in the source and
the first-line indentation LaTeX adds automatically to mark the paragraph break.

LaTeX allows basic formatting:

\textbf{Bold text}

\textit{Italic text}

\emph{Emphasised text}

Some characters have a special meaning in LaTeX and cannot simply be typed as-is.
The ten reserved characters are
\verb|# $ % & _ { } ~ ^ \|. To print seven of them literally, put a backslash
in front: \verb|\# \$ \% \& \_ \{ \}|. The remaining three are trickier: use
\verb|\textasciitilde| for \textasciitilde, \verb|\textasciicircum| for
\textasciicircum, and \verb|\textbackslash| for \textbackslash{} (writing
\verb|\\| would instead force a line break). Anything after an unescaped \%
becomes a comment, which is handy for leaving notes in your source but means a
value such as a 90\% win rate must be written \verb|90\%|. The tilde has a second
use outside \verb|\textasciitilde|: on its own, \verb|~| is a non-breaking space,
explained in Section~\ref{sec:sections}.

When you need to show code or literal symbols without any LaTeX interpretation,
use \verb|\verb| for a short inline snippet. It works by picking a delimiter
character that does not appear in your snippet and placing a copy on each side;
the vertical bar is used throughout this template, but any unused symbol works
(for example \verb+\verb+ with plus signs was used to typeset the bars you see
here). For whole blocks, such as a snippet of game code see Section~\ref{sec:code}.

A footnote adds a small note at the bottom of the page.\footnote{Like this one
--- useful for an aside or a source that is not a formal reference.} Write it
with \verb|\footnote{...}| placed \emph{immediately} after the word it attaches
to, with no space before it, so the raised marker sits tight against that word.
LaTeX numbers footnotes automatically.

Web addresses should be set with \verb|\url| (from the \texttt{url} package),
which uses a monospace font and lets a long link break sensibly across lines:
for example, the SDU site is \url{https://www.sdu.dk}. Do not type a bare URL as
ordinary text, because its \verb|~|, \verb|%|, or \verb|#| characters are special
(see above) and would need escaping. In a report, a footnote and a URL often go
together --- cite a web resource in a footnote using \verb|\footnote{\url{...}}|.

\begin{figure}[t]
    \centering
    \includegraphics[width=0.8\linewidth]{example-image}
    \caption{A screenshot of the prototype's first playable level.}
    \label{fig:single}
\end{figure}

% --------------------------------------------------
% Citations and References
% --------------------------------------------------

\section{Citations and References}
\label{sec:cite}

References are managed using a separate BibTeX file (e.g.,
\texttt{references.bib}). Each source is one entry in that file: an entry type
(such as \verb|@book|, \verb|@article|, or \verb|@inproceedings|), a citation
\emph{key} you choose, and a set of fields describing the source. For example,
the entry below defines a book with the key \texttt{tekinbas2003rules}:

\begin{verbatim}
@book{tekinbas2003rules,
  title={Rules of play: Game design
         fundamentals},
  author={Tekinbas, Katie Salen and
          Zimmerman, Eric},
  year={2003},
  publisher={MIT press}
}
\end{verbatim}

To cite that source in your text, use the \verb|\cite{}| command with its key:
writing \verb|\cite{tekinbas2003rules}| produces \cite{tekinbas2003rules}. You
never type the number yourself --- BibTeX assigns and orders the numbers, formats
each entry in the chosen style, and lists only the sources you actually cite.

A single reference can be cited as follows:
\cite{rigby2011immersion}

Multiple references can be cited together:
\cite{rigby2011immersion,salen2005game}

When several sources support the same statement, they can be grouped into one citation.
For example, research on player motivation and game design has been discussed extensively
in the literature \cite{rigby2011immersion,salen2005game,hunicke2004mda}.

The bibliography is generated automatically at the end of the document.
\verb|\bibliographystyle| selects the formatting (here, the IEEE style) and
\verb|\bibliography| names your \texttt{.bib} file without the extension:

\begin{verbatim}
\bibliographystyle{IEEEtran}
\bibliography{references}
\end{verbatim}

Note that BibTeX requires an extra build step: compile the document, run BibTeX,
then compile twice more so the citation numbers settle. Most editors (Overleaf,
TeXstudio, and others) do this for you automatically.

\subsection{Quotation Marks and Quotes}

LaTeX distinguishes opening from closing quotation marks, and the straight
typewriter character \verb|"| produces neither correctly. To get properly curved
marks, type two backticks to open and two apostrophes to close:
\verb|``...''| gives ``like this'', whereas \verb|"..."| gives the wrong shape.
The same applies to single quotes with \verb|`...'|.

For a short quotation, keep it inline within your own sentence and cite the
source. For a longer passage, set it apart using the \verb|quote| environment,
which indents the block so the reader can see at a glance that the words are not
yours. Always attribute a block quote with a \verb|\cite|:

\begin{quote}
Game design and authorship happen at many levels, and the fields of games
research and development involve people from diverse creative and scholarly
backgrounds. While it's often necessary to focus on one area, everyone,
regardless of discipline, will at some point need to consider issues outside
that area: base mechanisms of game systems, the overarching design goals, or the
desired experiential results of gameplay \cite{hunicke2004mda}.
\end{quote}

Use block quotes sparingly --- a report should mostly be your own analysis, with
quotations reserved for wording that matters.

% --------------------------------------------------
% Sections and Subsections
% --------------------------------------------------

\section{Sections and Subsections}
\label{sec:sections}

Documents are organised into a hierarchy of headings. \verb|\section| creates a
top-level heading (numbered automatically), \verb|\subsection| a level below it,
and \verb|\subsubsection| a level below that. You do not number them yourself;
LaTeX handles the numbering and will renumber everything if you reorder or insert
sections. You might use this structure to break your game down from a high-level
concept into individual systems.

To refer to a heading elsewhere, give it a \verb|\label| immediately after the
heading, then use \verb|\ref| to print its number. For example, this section was
labelled with \verb|\label{sec:sections}| and can be cited as
Section~\ref{sec:sections}; the citation guidance is in
Section~\ref{sec:cite}. Because LaTeX fills in the number, your cross-references
stay correct even if you reorder the document. A useful convention is to prefix
labels by type --- \texttt{sec:}, \texttt{fig:}, \texttt{tab:}, \texttt{lst:},
\texttt{eq:} --- so their purpose is obvious.

Notice the \verb|~| in \verb|Section~\ref{sec:sections}|. That tilde is a
\emph{non-breaking space}: it inserts an ordinary space but forbids a line break
at that point, so the word ``Section'' can never be separated from its number
across two lines. Always use \verb|~| between a label word and its reference
(\verb|Figure~\ref{...}|, \verb|Table~\ref{...}|), before a \verb|\cite|, and in
similar tight pairs such as ``Section~III'' or ``Elev~En''.


\subsection{Core Mechanics}

A subsection covers a more specific topic within its parent section, such as the
handful of core mechanics that define how your game plays.


\subsubsection{Enemy Behaviour}

A smaller heading can break a subsection down further --- for example, the
behaviour of a single enemy type within the game's combat system.


% --------------------------------------------------
% Lists
% --------------------------------------------------

\section{Lists}
\label{sec:lists}

Lists present a series of related items. LaTeX provides three built-in kinds,
each an environment whose entries are marked with \verb|\item|. Choose the kind
by what the order of the items means: use a bullet list when order does not
matter, a numbered list when it does (steps, rankings), and a description list
when each item is a term being defined.


\subsection{Bullet List}

A bullet list uses the \texttt{itemize} environment, for items with no
particular order, such as a game's headline features:

\begin{itemize}
    \item Local co-op multiplayer
    \item Procedurally generated levels
    \item Unlockable playable characters
\end{itemize}


\subsection{Numbered List}

A numbered list uses the \texttt{enumerate} environment, for ordered items such
as the steps of a design iteration. The numbers are generated automatically:

\begin{enumerate}
    \item Prototype the core mechanic
    \item Playtest with target players
    \item Iterate on the feedback
\end{enumerate}


\subsection{Description List}

A description list uses the \texttt{description} environment. The bracketed text
after each \verb|\item| becomes the bold term being defined, which is useful for
a glossary of game terms:

\begin{description}
    \item[Sprite] A 2D image drawn as part of the game world.
    \item[Hitbox] The region used to detect collisions.
    \item[Tick] A single fixed update of the game simulation.
\end{description}

Lists can also be nested inside one another, and the \texttt{enumitem} package
(already loaded) lets you customise the labels and spacing if needed.



% --------------------------------------------------
% Figures
% --------------------------------------------------

\section{Figures}
\label{sec:figures}

Images are placed with \verb|\includegraphics| inside a \texttt{figure}
environment, which turns the image into a numbered, captioned \emph{float} that
LaTeX positions for you. The optional argument in brackets is a \emph{placement
hint}: \verb|t| (top of a page), \verb|b| (bottom), \verb|h| (here,
approximately), and \verb|p| (a page of floats). You may combine letters --- the
hints are tried in the order LaTeX prefers, not the order you write them.

For this template, prefer \verb|[t]|: top-placed figures give the cleanest,
most consistent layout in a two-column document, and it is the IEEE convention.
If a figure will not sit where you want, or LaTeX pushes it far from its text,
relax the hint rather than fight it: \verb|[ht]| lets LaTeX try ``here'' before
falling back to the top, \verb|[h]| asks for roughly the current position, and
adding \verb|!| (as in \verb|[ht!]| or \verb|[!t]|) tells LaTeX to ignore its
usual spacing limits and honour your hint more aggressively. Reach for these
only when the default causes a real problem --- overusing \verb|h| tends to leave
awkward gaps. Set the image size relative to the text with
\verb|width=0.8\linewidth| rather than in fixed centimetres, so it scales with
the layout. Always give each figure a \verb|\caption| and a \verb|\label| so it
can be referenced; place the \verb|\label| after the \verb|\caption|.

Figures can then be referenced by number, e.g.\ Figure~\ref{fig:single} shows the
level. Because floats move, always refer to them by \verb|\ref| (with a
non-breaking \verb|~| before it) rather than by position (never ``the screenshot
below'').

A figure that needs the full width of the page, spanning both columns, uses the
starred form \verb|figure*| with \verb|width=\textwidth| --- useful for a wide
world map or a full UI mock-up. Note that two-column floats can only appear at
the top of a page, so they often move further from their source position than
single-column figures do.

\begin{figure*}[t]
  \centering
  \includegraphics[width=0.8\textwidth]{example-image-golden}
  \caption{The game's full world map, shown across both columns.}
  \label{fig:wide}
\end{figure*}

Figure~\ref{fig:wide} spans both columns in this way. 
This obviously has little to no effect in the one-column layout of the template.


% --------------------------------------------------
% Subfigures
% --------------------------------------------------

\section{Subfigures}
\label{sec:subfigures}

Several related images can be grouped inside a single figure using
\verb|\subfloat| (from the \texttt{subfig} package). Each subfigure gets its own
caption and its own \verb|\label|, so you can reference either the whole figure
or an individual part. The \verb|\hfill| between them spreads them evenly across
the column. This is ideal for showing how a level evolved between iterations.

\begin{figure}[t]
    \centering

    \subfloat[Greybox prototype]{
        \includegraphics[width=0.45\linewidth]{example-image-a}
        \label{fig:greybox}
    }
    \hfill
    \subfloat[Final art pass]{
        \includegraphics[width=0.45\linewidth]{example-image-b}
        \label{fig:finalart}
    }

    \caption{Two iterations of the same level: early greybox and final art.}
    \label{fig:iterations}
\end{figure}


Figure~\ref{fig:iterations} compares the two versions --- the greybox stage is
Figure~\ref{fig:greybox} and the final art is Figure~\ref{fig:finalart}. Note
that each subfigure label resolves to a letter (a), (b), while the parent figure
resolves to its number.


% --------------------------------------------------
% Tables
% --------------------------------------------------

\section{Tables}
\label{sec:tables}

Tables are useful for presenting structured information. The simplest form is a
few columns and rows inside a \verb|tabular| environment, wrapped in a
\verb|table| float so it can carry a caption and be referenced. The argument to
\verb|tabular| declares one column each: \verb|l|, \verb|c|, and \verb|r| give
left-, centre-, and right-aligned columns; rows are separated by \verb|\\|. In
professional IEEE style, tables use \emph{no vertical rules} and only a few
horizontal ones. The \texttt{booktabs} package provides these: \verb|\toprule|
at the top, \verb|\midrule| under the header, and \verb|\bottomrule| at the
bottom --- and its manual advises against ever adding vertical lines. Note also
that a table's \verb|\caption| goes \emph{above} it (unlike a figure's, which
goes below). Table~\ref{tab:controls} shows a basic example.

\begin{table}[!t]
\caption{Default Keyboard Controls}
\label{tab:controls}
\centering
\begin{tabular}{ll}
\toprule
\textbf{Action} & \textbf{Key}\\
\midrule
Move       & W A S D\\
Jump       & Space\\
Interact   & E\\
Pause      & Esc\\
\bottomrule
\end{tabular}
\end{table}

For richer data, a table can group columns under a shared heading with
\verb|\multicolumn| and let a label span several rows with \verb|\multirow|. A
\verb|\cmidrule| draws a partial rule under just the grouped columns.
Table~\ref{tab:playtest} summarises playtest results for two enemy AI variants
across two difficulty settings.

\begin{table}[!t]
\caption{Playtest Results by AI Variant and Difficulty}
\label{tab:playtest}
\centering
\begin{tabular}{llccr}
\toprule
\multirow{2}{*}{\textbf{AI Variant}} & \multirow{2}{*}{\textbf{Difficulty}}
  & \multicolumn{2}{c}{\textbf{Player Outcome}} & \multirow{2}{*}{\textbf{Rating}\textsuperscript{a}}\\
\cmidrule(lr){3-4}
 & & \textbf{Win \%} & \textbf{Avg.\ Time} & \\
\midrule
\multirow{2}{*}{Patrol}  & Normal & 72 & 4:15 & 3.8\\
                         & Hard   & 48 & 6:02 & 4.1\\
\midrule
\multirow{2}{*}{Flanker} & Normal & 55 & 5:30 & 4.4\\
                         & Hard   & 31 & 8:47 & 4.6\\
\bottomrule
\end{tabular}

\vspace{2pt}
{\footnotesize \textsuperscript{a}Mean fun rating on a 1--5 scale ($n=20$).}
\end{table}

Table~\ref{tab:playtest} presents the results. Note how
\verb|\multicolumn{2}{c}{...}| merges two header cells into one, and
\verb|\cmidrule(lr){3-4}| rules only the two columns under that grouped heading
(the \verb|(lr)| trims the rule slightly so it does not touch neighbouring
columns). A \verb|\midrule| separates the two AI variants. For a table too wide
for one column, wrap it in a \verb|table*| float so it spans both columns,
exactly as with a full-width figure.


% --------------------------------------------------
% Code Listings
% --------------------------------------------------

\section{Code Listings}\label{sec:code}

Source code is presented with the \texttt{lstlisting} environment, which
preserves indentation and adds syntax highlighting. Placing it inside a
\texttt{figure} gives it a caption, a label, and the same spacing as any other
figure. The listing below shows a \emph{fixed-timestep game loop}: physics is
advanced in constant-size steps so that simulation behaviour is independent of
frame rate, while rendering is interpolated to stay smooth on faster displays.

\begin{figure}[t]
\centering
\begin{lstlisting}
// Physics advances at a constant rate; rendering runs
// as fast as the platform allows.
const double dt = 1.0 / 60.0;   // 60 physics updates / second
double accumulator = 0.0;
double previous = getCurrentTime();

while (running) {
    double now = getCurrentTime();
    accumulator += now - previous;
    previous = now;

    processInput();

    // Consume fixed-size steps until we have caught up.
    while (accumulator >= dt) {
        updatePhysics(dt);
        accumulator -= dt;
    }

    // Blend factor for interpolating between physics states.
    double alpha = accumulator / dt;
    render(alpha);
}
\end{lstlisting}
\caption{Fixed-timestep game loop.}
\label{lst:gameloop}
\end{figure}

Listing~\ref{lst:gameloop} decouples the simulation rate from the render rate,
a common requirement for deterministic gameplay and networked physics. Change
the \verb|language=| key in the \verb|\lstset| block near the top of the file to
highlight a different language, such as C\# for Unity.


% --------------------------------------------------
% Algorithms
% --------------------------------------------------

\section{Algorithms}
\label{sec:algorithms}

Higher-level logic is often clearer as pseudocode than as source code. The
\texttt{algorithm} and \texttt{algorithmic} environments are used for this: the
\texttt{algorithm} wrapper provides the caption and number (like a float), while
\texttt{algorithmic} formats the steps. Commands such as \verb|\STATE|,
\verb|\IF|/\verb|\ENDIF|, \verb|\WHILE|/\verb|\ENDWHILE|, and \verb|\FORALL| set
the keywords in the conventional style, and the optional \verb|[1]| numbers the
lines. Algorithm~\ref{alg:astar} describes \emph{A* pathfinding} over a
navigation grid, a standard technique for moving enemies around obstacles.

\begin{algorithm}[t]
\caption{A* Pathfinding over a Navigation Grid}\label{alg:astar}
\begin{algorithmic}[1]
\REQUIRE start node $s$, goal node $t$, admissible heuristic $h(n)$
\ENSURE a shortest path from $s$ to $t$, or none if $t$ is unreachable
\STATE $open \gets \{s\}$ \COMMENT{priority queue ordered by $fScore$}
\STATE $gScore[s] \gets 0$
\STATE $fScore[s] \gets h(s)$
\WHILE{$open \neq \emptyset$}
    \STATE $current \gets \arg\min_{n \in open} fScore[n]$
    \IF{$current = t$}
        \RETURN \textsc{Reconstruct}($cameFrom, current$)
    \ENDIF
    \STATE remove $current$ from $open$
    \FORALL{neighbours $n$ of $current$}
        \STATE $tentative \gets gScore[current] + \textsc{Cost}(current, n)$
        \IF{$tentative < gScore[n]$}
            \STATE $cameFrom[n] \gets current$
            \STATE $gScore[n] \gets tentative$
            \STATE $fScore[n] \gets tentative + h(n)$
            \IF{$n \notin open$}
                \STATE add $n$ to $open$
            \ENDIF
        \ENDIF
    \ENDFOR
\ENDWHILE
\RETURN none
\end{algorithmic}
\end{algorithm}

The heuristic $h(n)$ estimates the remaining cost to the goal (for a tile grid,
Manhattan or Euclidean distance is typical). As long as $h$ never overestimates,
A* is guaranteed to return a shortest path for the enemy to follow.


% --------------------------------------------------
% Mathematics
% --------------------------------------------------

\section{Mathematics}
\label{sec:math}

Mathematics is written in \emph{math mode}. For a short formula inside a
sentence, use inline math delimited by dollar signs: the damage after a
$25\%$ armour reduction is $d' = 0.75\,d$. For a formula you want set apart on
its own line and automatically numbered, use the \texttt{equation} environment:

\begin{equation}
d' = d \left(1 - \frac{a}{a + k}\right)
\label{eq:armour}
\end{equation}

Equation~\ref{eq:armour} is a common armour-damage formula, where $a$ is armour
and $k$ a tuning constant, and \verb|\label| lets you refer back to it by number.
If you want a displayed formula \emph{without} a number, use \verb|\[ ... \]| (or
the \texttt{equation*} environment) instead.

Superscripts are written with \verb|^| and subscripts with \verb|_|. Each applies
to the single character that follows, so wrap anything longer in braces:
\verb|2^n| gives $2^n$ but \verb|2^{n-1}| is needed for $2^{n-1}$, and
\verb|v_{max}| gives $v_{max}$. The two can be combined on the same symbol, as in
\verb|x_i^2| for $x_i^2$. This is how you write things like a player's health
$h_i$ on frame $i$, or a score that grows as $2^{n}$ with the combo count $n$.

Beyond these, a few constructs come up constantly. Fractions use \verb|\frac|,
Greek letters have named commands such as \verb|\alpha| and \verb|\theta|, and
operators like $\sum$, $\sqrt{\;}$, and $\leq$ have their own commands too. For
instance, interpolating a character between two waypoints and averaging a score
over $n$ playtests can be written as Equation~\ref{eq:lerp}:

\begin{equation}
p(t) = (1-t)\,p_0 + t\,p_1, \qquad
\bar{s} = \frac{1}{n}\sum_{i=1}^{n} s_i
\label{eq:lerp}
\end{equation}

Note that spaces in the source are ignored in math mode, so LaTeX handles the
spacing for you; use \verb|\,| or \verb|\quad| only when you want to fine-tune
it. Remember also that symbols such as \verb|^|, \verb|_|, and Greek letters work
\emph{only} inside math mode --- writing them in ordinary text will produce an
error or the wrong output.


% --------------------------------------------------
% Conclusion
% --------------------------------------------------

\section{Conclusion}
\label{sec:conclusion}

Summarise the main outcomes of your game project: what you built, how it played,
and what your playtests showed. Keep it short --- the detail belongs in the
sections above.


% --------------------------------------------------
% Acknowledgments
% --------------------------------------------------

\section*{Acknowledgments}
The main use of the Acknowledgments section is to declare the usage of generative AI
in the project (e.g. assets, code) and the writing of the report.


% --------------------------------------------------
% References
% --------------------------------------------------

% \nocite{*}  % <= Uncomment this to show the whole biblography
\bibliographystyle{IEEEtran}
\bibliography{references}


% --------------------------------------------------
% Appendix
% --------------------------------------------------

\clearpage

\appendices

\section{Compiling This Template}
\label{app:setup}
This appendix explains how to build the template on Windows using \textbf{VS
Code} and \textbf{MiKTeX}. Follow the steps in order.

\begin{enumerate}
\item Download MiKTeX from \url{https://miktex.org/download} and run the
      installer. You can accept the defaults.
\item When asked to install missing packages on-the-fly, choose \textbf{Yes}.
      This lets MiKTeX fetch packages such as \texttt{booktabs} automatically.
\item Open the MiKTeX Console once and apply any updates it offers.
\item Download VS Code from \url{https://code.visualstudio.com} and install it.
\item In VS Code, open the Extensions view (\texttt{Ctrl+Shift+X}), search for
      \textbf{LaTeX Workshop}, and install it.
\item Open the Command Palette (\texttt{Ctrl+Shift+P}) and run
      \emph{Preferences: Open User Settings (JSON)}.
\item Type \emph{latex recipe} in the search box. 
      Set \emph{Latex-workshop Latex: Recipe.default} to
      \texttt{pdflatex -> bibtex -> pdflatex * 2}.
\item Unzip the template with \texttt{main.tex} and \texttt{references.bib} in one folder, 
      and open that folder with \emph{File $\rightarrow$ Open Folder}.
\item Open \texttt{main.tex} and save it (\texttt{Ctrl+S}). The PDF builds
      automatically; you can also press \texttt{Ctrl+Alt+B}.
\item The first build is slow. If MiKTeX asks to install packages, click
      \textbf{Install}. This happens only once.
\item Press \texttt{Ctrl+Alt+V} to view the PDF.
\end{enumerate}

If you get errors regarding pdflatex, Open the Command Palette 
(\texttt{Ctrl+Shift+P}) and run \emph{Preferences: Open User Settings (JSON)}.
Type \emph{latex recipe} in the search box. Click on \texttt{Edit in settings.json})
and paste the following block; overriding the previous setup.

\begin{footnotesize}
\begin{verbatim}
{
    "workbench.secondarySideBar.defaultVisibility": "hidden",
    "latex-workshop.latex.tools": [
    {
        "name": "pdflatex",
        "command": "pdflatex",
        "args": [
        "-synctex=1",
        "-interaction=nonstopmode",
        "-file-line-error",
        "%DOC%"
        ]
    },
    {
        "name": "bibtex",
        "command": "bibtex",
        "args": ["%DOCFILE%"]
    }
    ],
    "latex-workshop.latex.recipes": [
        
    {
        "name": "pdflatex -> bibtex -> pdflatex x2",
        "tools": ["pdflatex", "bibtex", "pdflatex", "pdflatex"]
    }
    ],
    "latex-workshop.latex.recipe.default": "first"
}
\end{verbatim}
\end{footnotesize}

If citations show as ``[?]'', build once more. If a build error names a missing
package, let MiKTeX install it and build again. If nothing happens on save,
check that you opened the folder (not just the file) and that LaTeX Workshop is
enabled.

\section{Full Playtest Questionnaire}
\label{app:questionnaire}
Appendices hold supporting material that would interrupt the main text --- raw
data, long tables, or full questionnaires. For a single appendix use
\verb|\appendix|; for several, use \verb|\appendices| and begin each one with
\verb|\section|, which LaTeX labels ``Appendix A'', ``Appendix B'', and so on.
Appendices can be labelled and referenced just like sections, so this one is
Appendix~\ref{app:questionnaire}. It would contain the complete list of questions
given to playtesters after each session, including the fun-rating scale
referenced in Table~\ref{tab:playtest}.

\section{Extended Controls Reference}
\label{app:controls}
This appendix (Appendix~\ref{app:controls}) would list the full keyboard and
gamepad control scheme, including the secondary and debug bindings omitted from
the summary in Table~\ref{tab:controls}.

\end{document}